\scanpage{44}
\mspage{35}

\[
\deleted{|Y|=\bigcup_{Z\in\widetilde Y}|Z|^{\circ}=\bigcup_{Z\in L}}
\]
alors les \deleted{$\unsure{X}$} $|\unsure{T|Z}|^{\circ}$ $(T\in M)$ sont \inserted{non vides et} deux à deux disjoints, \deleted{et on} et on a
\[
|Y|=\bigcup_{Z\in\widetilde Y}|T|^{\circ}
=\left(\bigcup_{Z\in L}|T|^{\circ}\right)\cup
\left(\bigcup_{Z\in\widetilde Y\setminus L}|T|^{\circ}\right)
\]
% margin note written slanted
\margin{ici on utilise (ii bis) non seulement pour $X$, mais aussi pour $Y$ et $Z$}
\deleted{$\unsure{R=} $}
et de même
\[
|Z|=\left(\bigcup_{Z\in L}|Z|^{\circ}\right)\cup
\left(\bigcup_{Z\in\unsure{\widetilde Y|\widetilde Z}\setminus L}|T|^{\circ}\right)
\]
d'où
\[
|Y|\cap|Z|=\unsure{-}\bigcup_{Z\in L}|Z|^{\circ},
\]
et ce n'est égal à $|Y|$, i.e. $|Y|\subset|Z|$, que si le deuxième \deleted{\unsure{est}} $\bigcup_{Z\in\widetilde Y\setminus L}|Z|^{\circ}\ne\varnothing$ i.e. $\widetilde Y\setminus L=\varnothing$, i.e. $Y\leq Z$; \deleted{\unsure{et}} on gagne ! Ainsi, on a que (ii bis) \inserted{(pour $X$ et les $Y\in\widetilde X$)} implique
\[
|X|=\bigcup_{Y\in\widetilde X}|\textcircled{X}\inserted{Y}|^{\circ}\tag*{(ii ter)}
\]
\margin{$\unsure{(\star)}\longrightarrow$}
\margin{\deleted{\unsure{c'est une}} c'est une partition indexée par $\widetilde X$, par (ii) et (iii), impliquant que $Y\ne Z$ $(Y,Z\in\widetilde X)\Rightarrow |Y|^{\circ}\cap|Z|^{\circ}=\varnothing$ donc $|Y|^{\circ}\ne|Z|^{\circ}$ (puisque $|Y|^{\circ}\ne\varnothing$ par (ii))}

Inversement, si on suppose (ii ter), a-t-on (ii bis) ?
\deleted{Voyons CF V p.~1}
\deleted{Pour la condition $3^{\circ}$ $(\alpha,\beta,\gamma)$. La condition $\unsure{\delta|\gamma}$ déjà vue \unsure{la première} (NB l'identité des $|X|^{\circ}$ définis ici, et de celui déduit de la}
Il faut voir (i) les $|X|^{\circ}\ne\varnothing$ (NB l'identité des $|X|^{\circ}$ définis ici, et de celui défini par la famille des $|\deleted{X}Y|$ $(Y\in\widetilde X)$ est maintenant immédiate \inserted{(cf \unsure{des lignes plus haut})}) c'est déjà vu, et (ii) $\forall x\in|X|$, $\exists$ un plus petit $Y\in\widetilde X$ avec $x\in|Y|^{\deleted{\unsure{\circ}}}$. Que la relation d'ordre de $\widetilde X$ soit celle induite par l'inclusion dans $\unsure{\mathscr P|\mathcal P}(P)$, a été vu par l'argument
% archive sheet number: 44
